\section{UART, Debug, and System Control}

\subsection{Why UART is useful}

A UART gives the CPU a simple way to print text, receive commands, and show program progress without a full display or operating system. Bare-metal programs often use UART as their first debug interface.

The repository contains several UART-related modules:

\begin{lstlisting}[language={}]
uart.v
uart_hw.v
uart_full.v
uart_rx.v
uart_tx.v
uart_tx_cfg.v
\end{lstlisting}

These are used by FPGA/SoC wrappers to provide serial I/O.

\subsection{Debug support}

\texttt{debug\_module.v}, \texttt{cpu\_core\_dbg.v}, and \texttt{soc\_dbg.v} are debug-oriented blocks/wrappers. The idea of a debug module is to let an external controller observe or influence CPU execution, for example halting, inspecting state, or stepping through code.

This appears experimental/educational rather than a full RISC-V Debug Specification implementation.

\subsection{System control}

\texttt{syscon.v} provides system-level control. In some SoC variants, a halt signal can freeze the CPU. This is useful in simulations, tests, or simple host-controlled systems.

\bigskip
